\section{Lineage and Context}
\label{sec:sources-wasm4pm-lineage}

\subsection{2016 Language-Model Lesson}
\label{sec:sources-wasm4pm-lineage-2016}

The dissertation's lineage originates in a 2016 language-model experiment in which local text
imitation was observed not to conserve consequence. A model could produce text that \emph{looked}
like an outcome — surface structure, vocabulary, even argument shape — without any mechanism
ensuring that the output \emph{followed} from the premises. That observation is the root of the
process-evidence program: prediction without consequence-preservation is noise dressed as
signal. The lesson applies directly to process mining. A fitness score of 0.97 produced by
a Python function is indistinguishable, to any downstream consumer, from a score that was
simply written into a field. Without a cryptographic receipt binding the score to the algorithm's
execution over a specific log, the score is \emph{text imitation of a measurement}, not a
measurement. \texttt{wasm4pm}'s receipt chain architecture is the direct technical response to
that 2016 lesson: every algorithm output must be bound to its provenance or it does not count.

\subsection{Enterprise Process Gap}
\label{sec:sources-wasm4pm-lineage-enterprise}

The enterprise process gap is the observation that organizations deploy process-automation
systems without a formal mechanism to verify that the declared process is the executed process.
The gap was apparent from process-intelligence research showing that real event logs, when mined,
reveal variant explosion, hidden loops, retries, and rework that the declared process model never
acknowledged. The Van der Aalst Constitution (\texttt{\textasciitilde/.claude/rules/process-mining-chicago-tdd.md})
formulates this as a hostile assumption: ``The declared manufacturing pipeline is not the real
runtime process.'' \texttt{wasm4pm} addresses the enterprise process gap by making conformance
checking an admission-gate mechanism rather than a reporting tool: a trace that does not meet the
fitness threshold is refused with a typed refusal, not merely flagged in a report. The
\texttt{ConformanceVerdict} enum (\texttt{FullyConforming} / \texttt{PartiallyConforming} /
\texttt{NonConforming}) and the 0.8 admission threshold in \texttt{src/conformance.rs} are the
direct technical instantiation of that enterprise-gap remediation.

\subsection{Relationship to the Chatman Equation}
\label{sec:sources-wasm4pm-lineage-chatman-equation}

The Chatman Equation \(A = \mu(O^*)\) asserts that enterprise value is a function of the
orchestrated process: value accrues when a correctly-specified process executes conformingly, and
erodes when the executed process deviates from the declared one. \texttt{wasm4pm} is the
executable substrate that makes $\mu$ computable. The Mining Authority discovers the process
model from the log (materializing $O^*$ from evidence); the Conformance Authority measures
$\mu$ as a fitness score; the Lifecycle Authority governs transitions across the six proof gates
that condition advancement on $\mu$ exceeding documented thresholds. Without \texttt{wasm4pm} as
the execution layer, the Chatman Equation remains a conceptual claim. With it, every evaluation
of $\mu$ is cryptographically receipted, append-only, and verifiable by a spot-audit framework.
[INTERPRETATION]

\subsection{Relationship to Receipts and Replay}
\label{sec:sources-wasm4pm-lineage-receipts-replay}

The receipt-and-replay strand of the dissertation lineage holds that every meaningful computation
must produce a cryptographic receipt that binds the output to the input, and that the receipt must
be replayable — a verifier must be able to re-run the computation and confirm the receipt. In
\texttt{wasm4pm}, this strand manifests at two levels. At the token level, the BLAKE3 receipt
implementation (\texttt{src/crypto.rs}, self-implemented with Merkle-tree chaining, zero external
dependencies) computes rolling \texttt{Blake3Hash([u8;~32])} fields that bind each
\texttt{Evidence<T,~S,~W>} container to the prior chain state. At the algorithm level, the
\texttt{TokenReplayEngine} (van der Aalst fitness equation $f = 0.5(1-m/c) + 0.5(1-r/p)$ with
source/sink token accounting) and the \texttt{AlignmentEngine} (A* search in synchronous product
space, 5000-iteration cap, \texttt{BinaryHeap}) both emit \texttt{ConformanceVerdicts} that carry
aggregate fitness and per-case verdict histories. The open gap (documented as GAP\_009) is that
BLAKE3 receipt fields in \texttt{Evidence} are currently initialized to all-zero hashes: the
infrastructure is present but the hash computation over algorithm inputs is not yet wired end-to-end.
